\chapter{A Dignidade da Pessoa Humana}

O terceiro grande bloco do Catecismo da Igreja Católica abre com uma questão fundante: o que é o homem e para que existe? A resposta da fé cristã não provém de especulação filosófica, mas da revelação de Cristo, que ``ao revelar plenamente o mistério do Pai e do seu amor, manifesta plenamente o homem a si mesmo e descobre-lhe a sua vocação sublime'' (GS 22). Nos parágrafos 1700 a 1876, o Catecismo percorre os fundamentos da dignidade humana: a criação à imagem de Deus, a vocação à bem-aventurança, a liberdade, a moralidade dos atos, as paixões, a consciência, as virtudes e o pecado. Todo este percurso é orientado por uma convicção central: o homem está chamado a \textbf{fazer o bem} --- a agir de modo a conduzir a sua vida ao seu fim último, que é o próprio Deus.

\section{O homem, imagem de Deus (1700--1715)}

A dignidade da pessoa humana tem a sua raiz na criação à imagem e semelhança de Deus. Foi em Cristo, ``imagem do Deus invisível'' (Cl 1,15), que o homem foi criado à imagem do Criador; e foi em Cristo, redentor e salvador, que essa imagem divina, deformada no homem pelo primeiro pecado, foi restaurada na sua beleza original e enobrecida pela graça de Deus. Esta imagem divina está presente em cada homem e resplandece na comunhão das pessoas, à semelhança da unidade das Pessoas divinas entre si. Toda a dignidade do homem parte, portanto, desta verdade fundamental, que não é um dado apenas filosófico mas o núcleo da revelação cristã.

Dotada de uma alma espiritual e imortal, a pessoa humana é ``a única criatura sobre a terra querida por Deus por si mesma'' (GS 24). Desde que é concebida, é destinada à bem-aventurança eterna. Pela razão, é capaz de compreender a ordem das coisas estabelecida pelo Criador; pela vontade, é capaz de se orientar para o \textbf{bem} verdadeiro. E encontra a sua perfeição na ``busca e no amor da verdade e do \textbf{bem}'' (GS 15). Em virtude das forças espirituais da inteligência e da vontade, o homem é ainda dotado de liberdade, ``sinal privilegiado da imagem divina'' (GS 17). Mediante a razão, conhece a voz de Deus que o impele a ``\textbf{fazer o bem} e a evitar o mal'' (GS 16): todos devem seguir esta lei moral, que ressoa na consciência e se cumpre no amor de Deus e do próximo.

Todavia, ``seduzido pelo Maligno desde o começo da história, o homem abusou da sua liberdade'' (GS 13). Sucumbiu à tentação e cometeu o mal. Conserva o desejo do \textbf{bem}, mas a sua natureza está ferida pelo pecado original: ficou com a inclinação para o mal e sujeito ao erro, de tal forma que ``toda a vida humana, quer singular quer colectiva, se apresenta como uma luta, e quão dramática, entre o \textbf{bem} e o mal, entre a luz e as trevas'' (GS 13). Pela sua paixão, porém, Cristo livrou-nos de Satanás e do pecado e mereceu-nos a vida nova no Espírito Santo; a sua graça restaura o que o pecado havia deteriorado. Quem crê em Cristo torna-se filho de Deus: esta adoção filial capacita-o a seguir o exemplo de Cristo, a agir com retidão e a \textbf{fazer o bem}. Amadurecida na graça, a vida moral culmina na vida eterna, na glória do céu.

\section{Nossa vocação à bem-aventurança (1716--1729)}

As bem-aventuranças estão no coração da pregação de Jesus. O seu anúncio retoma as promessas feitas ao povo eleito desde Abraão e completa-as, ordenando-as, não já somente à felicidade resultante da posse de uma terra, mas ao Reino dos céus: ``Bem-aventurados os pobres em espírito, porque deles é o Reino dos céus \ldots'' (Mt 5,3--12). As bem-aventuranças retratam o rosto de Jesus Cristo e descrevem a sua caridade; exprimem a vocação dos fiéis associados à glória da sua paixão e ressurreição; são as promessas paradoxais que sustentam a esperança no meio das tribulações, anunciando já, de modo obscuro, as bênçãos e recompensas adquiridas pelos discípulos de Cristo. Respondem, assim, ao desejo natural de felicidade que Deus colocou no coração do homem para o atrair a Si, o único que o pode satisfazer: ``Só Deus sacia'' (São Tomás de Aquino).

As bem-aventuranças revelam a meta da existência humana, o fim último dos atos humanos: Deus chama-nos à sua própria felicidade. Esta vocação dirige-se a cada um, pessoalmente, mas também ao conjunto da Igreja, povo novo constituído por aqueles que acolheram a promessa e dela vivem na fé. O Novo Testamento emprega várias expressões para descrever esta bem-aventurança: a chegada do Reino de Deus, a visão de Deus --- ``Bem-aventurados os puros de coração, porque verão a Deus'' (Mt 5,8) ---, a entrada na alegria do Senhor, a entrada no repouso de Deus. Uma tal bem-aventurança ultrapassa a inteligência e as simples forças humanas; resulta de um dom gratuito de Deus, e por isso se classifica de sobrenatural. De facto, Deus colocou-nos no mundo para O conhecermos, servirmos e amarmos, e assim chegarmos ao paraíso. Com a bem-aventurança, o homem entra na glória de Cristo e no gozo da vida trinitária.

A bem-aventurança prometida coloca-nos perante as opções morais decisivas: convida-nos a purificar o coração dos seus maus instintos e a procurar o amor de Deus acima de tudo, ensinando que a verdadeira felicidade não reside na riqueza ou no bem-estar, na glória humana ou no poder, em nenhuma obra ou criatura, mas só em Deus, fonte de todo o \textbf{bem} e de todo o amor. O Decálogo, o Sermão da Montanha e a catequese apostólica descrevem os caminhos que conduzem ao Reino dos céus; por eles avançamos, passo a passo, pelos atos de cada dia, amparados pela graça do Espírito Santo. A bem-aventurança do céu determina os critérios de discernimento no uso dos bens terrenos, em conformidade com a Lei de Deus: a vocação à bem-aventurança não afasta o homem do mundo, mas orienta os seus atos para o fim último que lhes confere verdadeiro sentido.

\section{A liberdade do homem (1730--1748)}

Deus criou o homem racional, conferindo-lhe a dignidade de pessoa dotada de iniciativa e do domínio dos seus próprios atos. ``Deus quis `deixar o homem entregue à sua própria decisão' (Sir 15,14), de tal modo que procure por si mesmo o seu Criador e, aderindo livremente a Ele, chegue à total e beatífica perfeição.'' A liberdade é o poder, radicado na razão e na vontade, de agir ou não agir, de fazer isto ou aquilo. Pelo livre-arbítrio, cada qual dispõe de si: a liberdade é, no homem, ``uma força de crescimento e de maturação na verdade e na bondade'', atingindo a sua perfeição quando está ordenada para Deus, nossa bem-aventurança. Enquanto não se fixa definitivamente no seu \textbf{bem} último, que é Deus, a liberdade implica a possibilidade de escolher entre o \textbf{bem} e o mal, e portanto de crescer na perfeição ou de falhar e pecar.

Quanto mais o homem \textbf{fizer o bem}, mais livre se torna: não há verdadeira liberdade senão no serviço do \textbf{bem} e da justiça. A liberdade torna o homem responsável pelos seus atos, na medida em que são voluntários; o progresso na virtude, o conhecimento do \textbf{bem} e a ascese aumentam o domínio da vontade sobre os próprios atos. A imputabilidade e responsabilidade de um ato podem ser diminuídas ou até anuladas pela ignorância, a inadvertência, a violência, o medo, os hábitos e as afeições desordenadas. Mas o exercício da liberdade não implica o direito de tudo dizer e de tudo fazer. O homem que rejeita o projeto divino de amor engana-se a si mesmo e torna-se escravo do pecado; afastando-se da lei moral, atenta contra a sua própria liberdade, agrilhoa-se a si mesmo e rebela-se contra a verdade divina.

Pela sua cruz gloriosa, Cristo obteve a salvação de todos os homens, resgatando-os do pecado que os retinha em situação de escravatura. ``Foi para a liberdade que Cristo nos libertou'' (Gl 5,1): foi-nos dado o Espírito Santo e, como ensina o Apóstolo, ``onde está o Espírito, aí está a liberdade'' (2 Cor 3,17). A graça de Cristo não faz concorrência à nossa liberdade quando esta corresponde ao sentido da verdade e do \textbf{bem} que Deus colocou no coração do homem; pelo contrário, quanto mais dóceis formos aos impulsos da graça, tanto mais crescem a nossa liberdade interior e a nossa segurança nas provações e nos constrangimentos exteriores. Pela ação da graça, o Espírito Santo educa-nos para a liberdade espiritual, tornando-nos colaboradores livres da sua obra na Igreja e no mundo.

\section{A moralidade dos atos humanos (1749--1761)}

A liberdade faz do homem um sujeito moral. Quando age de maneira deliberada, o homem é, por assim dizer, o pai dos seus atos. Os atos humanos, livremente escolhidos em consequência de um juízo de consciência, são moralmente qualificáveis: são bons ou maus. A moralidade dos atos humanos depende do objeto escolhido, do fim que se tem em vista (a intenção) e das circunstâncias da ação --- estas três são as ``fontes'' ou elementos constitutivos da moralidade dos atos humanos. O objeto escolhido é um \textbf{bem} para o qual a vontade tende deliberadamente; as regras objetivas da moralidade enunciam a ordem racional do \textbf{bem} e do mal, atestada pela consciência. O objeto escolhido especifica moralmente o ato da vontade na medida em que a razão o reconhece e o julga conforme, ou não, ao verdadeiro \textbf{bem}.

A intenção coloca-se do lado do sujeito que age: é a orientação da vontade em direção ao fim, o alvo do \textbf{bem} que se espera da ação empreendida. Uma intenção boa não torna, porém, bom nem justo um comportamento em si mesmo desordenado: ``o fim não justifica os meios.'' Inversamente, uma intenção má acrescentada torna mau um ato que em si poderia ser bom. As circunstâncias, incluindo as consequências, são elementos secundários de um ato moral; podem agravar ou atenuar a bondade ou malícia moral dos atos, mas não podem de per si modificar a qualidade moral dos próprios atos: não podem tornar boa nem justa uma ação má em si mesma.

O ato moralmente \textbf{bom} pressupõe, em simultâneo, a bondade do objeto, da finalidade e das circunstâncias: um fim mau corrompe a ação mesmo que o seu objeto seja \textbf{bom} em si. Há comportamentos concretos cuja escolha é sempre um erro, porque comportam uma desordem da vontade, um mal moral: são atos que, por si e em si mesmos, independentemente das circunstâncias e das intenções, são sempre gravemente ilícitos em razão do seu objeto, como a blasfémia, o jurar falso, o homicídio e o adultério. ``Não é permitido \textbf{fazer o mal} para que daí resulte um \textbf{bem}.'' Esta verdade protege a integridade da vida moral e salvaguarda a pessoa de manipulações que instrumentalizem o \textbf{bem} aparente para justificar ações intrinsecamente más.

\section{A moralidade das paixões (1762--1775)}

A pessoa humana ordena-se à bem-aventurança através dos seus atos deliberados; as paixões ou sentimentos que experimenta podem dispô-la nesse sentido e contribuir para isso. O termo ``paixões'' pertence ao patrimônio cristão e designa as emoções ou movimentos da sensibilidade que inclinam a agir, ou a não agir, em vista do que se sentiu ou imaginou como \textbf{bom} ou como mau. As paixões são componentes naturais do psiquismo humano; constituem o lugar de passagem e garantem a ligação entre a vida sensível e a vida do espírito. Nosso Senhor designa o coração do homem como a fonte de onde brota o movimento das paixões (cf.~Mc 7,21).

São numerosas as paixões. A mais fundamental é o amor, provocado pela atração do \textbf{bem}: o amor causa o desejo do \textbf{bem} ausente e a esperança de o alcançar, e este movimento tem o seu termo no prazer e na alegria do \textbf{bem} possuído. A apreensão pelo mal, por outro lado, causa o ódio, a aversão e o receio do mal futuro, terminando na tristeza pelo mal presente ou na cólera que a ele se opõe. ``Amar é querer \textbf{bem} a alguém'' (São Tomás de Aquino). Todos os outros afetos nascem neste movimento original do coração para o \textbf{bem}: só o \textbf{bem} é amado, e por isso ``as paixões são más se o amor for mau, e boas se ele for bom'' (Santo Agostinho).

Em si mesmas, as paixões não são nem boas nem más; só recebem qualificação moral na medida em que dependem efetivamente da razão e da vontade. Pertence à perfeição do \textbf{bem} moral que as paixões sejam reguladas pela razão. A vontade reta ordena para o \textbf{bem} e para a bem-aventurança os movimentos sensíveis que assume; a vontade má sucumbe às paixões desordenadas e exacerba-as. As emoções e os sentimentos podem ser assumidos pelas virtudes ou pervertidos pelos vícios. Na vida cristã, o próprio Espírito Santo realiza a sua obra mobilizando todo o ser; em Cristo, os sentimentos humanos podem alcançar a sua consumação na caridade e na bem-aventurança divina. ``A perfeição moral consiste em que o homem não seja movido para o \textbf{bem} só pela vontade, mas também pelo apetite sensível'' (cf.~Sl 84,3): a integração das paixões na vida moral é sinal da maturidade do discípulo.

\section{A consciência moral (1776--1802)}

``No mais profundo da consciência, o homem descobre uma lei que não se deu a si mesmo, mas à qual deve obedecer e cuja voz ressoa, quando necessário, aos ouvidos do seu coração, chamando-o sempre a amar e \textbf{fazer o bem} e a evitar o mal [\ldots]. A consciência é o núcleo mais secreto e o sacrário do homem, no qual ele se encontra a sós com Deus, cuja voz ressoa na intimidade do seu ser'' (GS 16). Presente no coração da pessoa, a consciência moral leva-a, no momento oportuno, a \textbf{fazer o bem} e a evitar o mal; ela julga as opções concretas, aprovando as boas e denunciando as más, atestando a autoridade da verdade em relação ao \textbf{bem} supremo pelo qual a pessoa humana se sente atraída. A consciência moral é um juízo da razão, pelo qual a pessoa reconhece a qualidade moral de um ato concreto que vai praticar, que está prestes a executar ou que já realizou.

A dignidade da pessoa humana implica e exige a retidão da consciência moral. Esta compreende a perceção dos princípios da moralidade, a sua aplicação em determinadas circunstâncias por meio de um discernimento prático das razões e dos \textbf{bens}, e, por fim, o juízo emitido sobre os atos concretos. Se o homem comete o mal, o justo juízo da consciência pode ser nele a testemunha da verdade universal do \textbf{bem} e, ao mesmo tempo, da maldade da sua opção concreta; mas ``o veredicto do juízo da consciência continua a ser um penhor de esperança e de misericórdia'', lembrando o perdão a pedir, o \textbf{bem} a praticar ainda e a virtude a cultivar incessantemente com a graça de Deus. O homem tem também o direito de agir em consciência e em liberdade, não sendo forçado a agir contra ela, sobretudo em matéria religiosa.

A consciência deve ser informada e o juízo moral esclarecido: uma consciência bem formada é reta e verídica, formulando os seus juízos segundo a razão e em conformidade com o \textbf{bem} verdadeiro querido pela sabedoria do Criador. A formação da consciência é tarefa para toda a vida: a Palavra de Deus é a luz do nosso caminho, devemos assimilá-la na fé e na oração e pô-la em prática. Perante a necessidade de decidir moralmente, algumas regras se aplicam sempre: nunca é permitido \textbf{fazer o mal} para que daí resulte um \textbf{bem}; a ``regra de ouro'' é ``Tudo quanto quiserdes que os homens vos façam, fazei-lho, de igual modo, vós também'' (Mt 7,12); e a caridade passa sempre pelo respeito do próximo e da sua consciência. A consciência pode, porém, permanecer na ignorância ou fazer juízos errôneos; tal ignorância e tais erros nem sempre são isentos de culpabilidade, pois a ignorância voluntária e o endurecimento do coração não diminuem, antes aumentam, o caráter voluntário do pecado.

\section{As virtudes (1803--1845)}

``Tudo o que é verdadeiro, nobre e justo, tudo o que é puro, amável e de boa reputação, tudo o que é virtude e digno de louvor, isto deveis ter no pensamento'' (Fl 4,8). A virtude é uma disposição habitual e firme para praticar o \textbf{bem}: permite à pessoa não somente praticar atos \textbf{bons}, mas dar o melhor de si mesma, tendendo para o \textbf{bem} com todas as suas forças sensíveis e espirituais, procurando-o e optando por ele em atos concretos. ``O fim de uma vida virtuosa consiste em tornar-se semelhante a Deus'' (São Gregório de Nissa). As virtudes humanas são atitudes firmes, disposições estáveis, perfeições habituais da inteligência e da vontade, que regulam os nossos atos, ordenam as nossas paixões e guiam o nosso procedimento segundo a razão e a fé; conferem facilidade, domínio e alegria para se levar uma vida moralmente boa, sendo os frutos e os germes de atos moralmente \textbf{bons}.

Há quatro virtudes que desempenham um papel central e chamam-se ``cardeais''; todas as outras se agrupam em torno delas: a prudência, a justiça, a fortaleza e a temperança. A prudência é a ``recta norma da acção'' que dispõe a razão prática para discernir, em qualquer circunstância, o nosso verdadeiro \textbf{bem} e para escolher os justos meios de o atingir: é ela que guia imediatamente o juízo da consciência. A justiça consiste na constante e firme vontade de dar a Deus e ao próximo o que lhes é devido. A fortaleza assegura, no meio das dificuldades, a firmeza e a constância na prossecução do \textbf{bem}, capacitando a vencer o medo, enfrentar a provação e ir até à renúncia e ao sacrifício da própria vida na defesa de uma causa justa. A temperança modera a atração dos prazeres e proporciona o equilíbrio no uso dos bens criados, orientando para o \textbf{bem} os apetites sensíveis e guardando uma sã discrição. ``Viver bem é amar a Deus de todo o coração, com toda a alma e com todo o proceder'' (Santo Agostinho).

As virtudes humanas radicam nas virtudes teologais --- fé, esperança e caridade ---, que têm Deus por origem, motivo e objeto, e dispõem os cristãos para viverem em relação com a Santíssima Trindade. A fé é a virtude pela qual cremos em Deus e em tudo o que Ele nos revelou; o justo vive pela fé, que ``atua pela caridade'' (Gl 5,6). A esperança é a virtude pela qual desejamos o Reino dos céus e a vida eterna como nossa felicidade, pondo toda a confiança nas promessas de Cristo: sustenta no abatimento, dilata o coração na expectativa da bem-aventurança eterna e preserva do egoísmo, conduzindo à felicidade da caridade. A caridade é a virtude pela qual amamos a Deus sobre todas as coisas por Ele mesmo, e ao próximo como a nós mesmos por amor de Deus; é o ``vínculo da perfeição'' (Cl 3,14) e a forma de todas as virtudes. ``A prática da vida moral animada pela caridade dá ao cristão a liberdade espiritual dos filhos de Deus'': o cristão não está diante de Deus como escravo nem como mercenário, mas como filho que corresponde ao amor d'Aquele que nos amou primeiro. A vida moral dos cristãos é ainda sustentada pelos sete dons do Espírito Santo --- sabedoria, entendimento, conselho, fortaleza, ciência, piedade e temor de Deus ---, que tornam os fiéis dóceis, na obediência pronta, às inspirações divinas.

\section{O pecado (1846--1876)}

O Evangelho é a revelação, em Jesus Cristo, da misericórdia de Deus para com os pecadores (cf.~Lc 15). ``Deus, que nos criou sem nós, não quis salvar-nos sem nós'' (Santo Agostinho): o acolhimento da sua misericórdia exige de nós a confissão das nossas faltas. ``Onde abundou o pecado, superabundou a graça'' (Rm 5,20); mas para realizar a sua obra, a graça tem de pôr a descoberto o pecado, para converter o nosso coração. Como um médico que examina a chaga antes de lhe aplicar o penso, Deus, pela sua Palavra e pelo seu Espírito, projeta uma luz viva sobre o pecado, tornando possível a conversão que contém ``o juízo interior da consciência'' e torna-se ``o princípio de um novo dom da graça e do amor'' (João Paulo II, \textit{Dominum et vivificantem}, 31).

O pecado é uma falta contra a razão, a verdade, a reta consciência; é uma falha contra o verdadeiro amor para com Deus e para com o próximo, por causa de um apego perverso a certos bens. Fere a natureza do homem e atenta contra a solidariedade humana. É ``uma palavra, um ato ou um desejo contrários à lei eterna'' (Santo Agostinho). O pecado é uma ofensa a Deus; é, como o primeiro pecado, uma desobediência --- uma revolta contra Deus pela vontade de se tornarem ``como deuses'', conhecendo e determinando por si mesmos o que é \textbf{bem} e o que é mal (Gn 3,5). Assim, o pecado é ``o amor de si próprio levado até ao desprezo de Deus'' (Santo Agostinho), diametralmente oposto à obediência de Jesus, que realizou a salvação. Os pecados distinguem-se segundo o seu objeto e segundo a sua gravidade: o pecado mortal destrói a caridade no coração do homem por uma infração grave à Lei de Deus, desviando-o do seu último fim e da sua bem-aventurança; o pecado venial deixa subsistir a caridade, embora ofendendo-a e ferindo-a. Para que um pecado seja mortal, requerem-se em simultâneo: matéria grave, plena consciência e deliberado consentimento.

O pecado arrasta ao pecado: gera o vício, pela repetição dos mesmos atos, resultando em inclinações perversas que obscurecem a consciência e corrompem a apreciação concreta do \textbf{bem} e do mal. Os pecados capitais --- soberba, avareza, inveja, ira, luxúria, gula e preguiça --- são assim chamados porque são geradores de outros pecados e vícios. O pecado é ainda um ato pessoal com responsabilidade social: ``o pecado torna os homens cúmplices uns dos outros, faz reinar entre eles a concupiscência, a violência e a injustiça''; as ``estruturas de pecado'' são expressão e efeito dos pecados pessoais e induzem as suas vítimas a que, por sua vez, cometam o mal. Mas a última palavra não pertence ao pecado: ``Deus encerrou todos na desobediência, para usar de misericórdia para com todos'' (Rm 11,32), e a misericórdia de Deus, manifestada plenamente em Jesus Cristo crucificado e ressuscitado, é sempre maior que o pecado, abrindo para todo o homem o caminho do arrependimento, da reconciliação e do retorno ao \textbf{bem}.
